% GENERATED FILE. Edit canonical lessons, then run npm run generate:books.
% canonical-source-hash: fnv1a64:acc19ad1e12838fc
% canonical-lessons: SA-C490-gaurah, SA-C490-ayuh, SA-C490-palitah, SA-C490-valih, SA-C490-rupam

\chapter{Age and Appearance}
\label{ch:sa-age-and-appearance}
\hlchaptermodality{490}

\begin{chapteropening}
\textbf{By the end of this chapter:} \emph{I can describe someone's age and appearance.}

Complete the last lesson of chapter 490: I can describe someone's age and appearance.
\end{chapteropening}

\section[\texorpdfstring{\sk{गौरः} (gauraḥ)}{गौरः (gauraḥ)}]{\sk{गौरः} (gauraḥ) — pale; whitish}

\label{lesson:SA-C490-gaurah}



\begin{quote}
Before the new one: say the Sanskrit for gives up entirely, then the Sanskrit for liquor.
\end{quote}

\subsection*{You'll want to know: \sk{गौरः}}
\textbf{\sk{गौरः}} — \emph{gauraḥ} — \textquotedblleft{}pale; whitish\textquotedblright{}.

\begin{quote}
\sk{शिशुः} \sk{गौरः} \sk{सुन्दरः} \sk{च} \sk{अस्ति।}
\end{quote}

\begin{quote}
\emph{śiśuḥ gauraḥ sundaraḥ ca asti.}
\end{quote}

\begin{quote}
The baby is fair and lovely.
\end{quote}

\begin{cousinweb}
\textbf{\sk{गौरः}} (\emph{gauraḥ}) MW gives \textquotedblleft{}white, yellowish, reddish, pale red\textquotedblright{}; MW gives no derivation, so its origin is uncertain (MW).
\end{cousinweb}

\begin{grammarlens}[title={What you've built}]
One more word for describing what you lived and felt.
\end{grammarlens}

\subsection*{Your turn}
\begin{itemize}
\raggedright
  \item \emph{Say it:} \emph{gauraḥ}
  \item \emph{Say it:} \emph{gauraḥ}, once more
  \item \emph{Say it:} \emph{śiśuḥ gauraḥ sundaraḥ ca asti.} — The baby is fair and lovely.
  \item \emph{Recall:} say the Sanskrit for gives up entirely, then the Sanskrit for liquor, then say \emph{gauraḥ} again
\end{itemize}

\begin{tcolorbox}[breakable,colback=teal!4,colframe=teal!35!black,title={Before you move on}]
What does \sk{गौरः} mean? (\textquotedblleft{}Pale; whitish\textquotedblright{}.) Say it once more.
\end{tcolorbox}
\section[\texorpdfstring{\sk{आयुः} (āyuḥ)}{आयुः (āyuḥ)}]{\sk{आयुः} (āyuḥ) — length of life; life}

\label{lesson:SA-C490-ayuh}



\begin{quote}
Before the new one: say the Sanskrit for liquor, then the Sanskrit for pale.
\end{quote}

\subsection*{You'll want to know: \sk{आयुः}}
\textbf{\sk{आयुः}} — \emph{āyuḥ} — \textquotedblleft{}length of life; life\textquotedblright{}.

\begin{quote}
\sk{व्यायामेन} \sk{आयुः} \sk{वर्धते।}
\end{quote}

\begin{quote}
\emph{vyāyāmena āyuḥ vardhate.}
\end{quote}

\begin{quote}
Exercise lengthens life.
\end{quote}

A neuter noun.

\begin{cousinweb}
\textbf{\sk{आयुः}} (\emph{āyuḥ}) MW gives \textquotedblleft{}life, vital power, health, duration of life, long life\textquotedblright{}, the word at the head of \textbf{\sk{आयुर्वेदः}}; MW gives no derivation, so its origin is uncertain (MW).
\end{cousinweb}

\begin{grammarlens}[title={What you've built}]
One more word for describing what you lived and felt.
\end{grammarlens}

\subsection*{Your turn}
\begin{itemize}
\raggedright
  \item \emph{Say it:} \emph{āyuḥ}
  \item \emph{Say it:} \emph{āyuḥ}, once more
  \item \emph{Say it:} \emph{vyāyāmena āyuḥ vardhate.} — Exercise lengthens life.
  \item \emph{Recall:} say the Sanskrit for liquor, then the Sanskrit for pale, then say \emph{āyuḥ} again
\end{itemize}

\begin{tcolorbox}[breakable,colback=teal!4,colframe=teal!35!black,title={Before you move on}]
What does \sk{आयुः} mean? (\textquotedblleft{}Length of life; life\textquotedblright{}.) Say it once more.
\end{tcolorbox}
\section[\texorpdfstring{\sk{पलितः} (palitaḥ)}{पलितः (palitaḥ)}]{\sk{पलितः} (palitaḥ) — grey-haired}

\label{lesson:SA-C490-palitah}



\begin{quote}
Before the new one: say the Sanskrit for pale, then the Sanskrit for length of life.
\end{quote}

\subsection*{You'll want to know: \sk{पलितः}}
\textbf{\sk{पलितः}} — \emph{palitaḥ} — \textquotedblleft{}grey-haired\textquotedblright{}.

\begin{quote}
\sk{पितामहस्य} \sk{केशाः} \sk{पलिताः} \sk{सन्ति।}
\end{quote}

\begin{quote}
\emph{pitāmahasya keśāḥ palitāḥ santi.}
\end{quote}

\begin{quote}
Grandfather's hair is grey.
\end{quote}

\begin{cousinweb}
\textbf{\sk{पलितः}} (\emph{palitaḥ}) MW gives \textquotedblleft{}grey, hoary, old\textquotedblright{} and the noun \textquotedblleft{}grey hair\textquotedblright{}; MW gives no derivation, so its origin is uncertain (MW).
\end{cousinweb}

\begin{grammarlens}[title={What you've built}]
One more word for describing what you lived and felt.
\end{grammarlens}

\subsection*{Your turn}
\begin{itemize}
\raggedright
  \item \emph{Say it:} \emph{palitaḥ}
  \item \emph{Say it:} \emph{palitaḥ}, once more
  \item \emph{Say it:} \emph{pitāmahasya keśāḥ palitāḥ santi.} — Grandfather's hair is grey.
  \item \emph{Recall:} say the Sanskrit for pale, then the Sanskrit for length of life, then say \emph{palitaḥ} again
\end{itemize}

\begin{tcolorbox}[breakable,colback=teal!4,colframe=teal!35!black,title={Before you move on}]
What does \sk{पलितः} mean? (\textquotedblleft{}Grey-haired\textquotedblright{}.) Say it once more.
\end{tcolorbox}
\section[\texorpdfstring{\sk{वलिः} (valiḥ)}{वलिः (valiḥ)}]{\sk{वलिः} (valiḥ) — a wrinkle}

\label{lesson:SA-C490-valih}



\begin{quote}
Before the new one: say the Sanskrit for length of life, then the Sanskrit for grey-haired.
\end{quote}

\subsection*{You'll want to know: \sk{वलिः}}
\textbf{\sk{वलिः}} — \emph{valiḥ} — \textquotedblleft{}a wrinkle\textquotedblright{}.

\begin{quote}
\sk{वृद्धस्य} \sk{ललाटे} \sk{वलिः} \sk{अस्ति।}
\end{quote}

\begin{quote}
\emph{vṛddhasya lalāṭe valiḥ asti.}
\end{quote}

\begin{quote}
There is a wrinkle on the old man's forehead.
\end{quote}

A feminine noun.

\begin{cousinweb}
\textbf{\sk{वलिः}} (\emph{valiḥ}) MW gives \textquotedblleft{}a fold of the skin, a wrinkle\textquotedblright{} and lists the spellings \emph{bali} and \emph{valī} beside it; MW gives no derivation, so its origin is uncertain (MW).
\end{cousinweb}

\begin{grammarlens}[title={What you've built}]
One more word for describing what you lived and felt.
\end{grammarlens}

\subsection*{Your turn}
\begin{itemize}
\raggedright
  \item \emph{Say it:} \emph{valiḥ}
  \item \emph{Say it:} \emph{valiḥ}, once more
  \item \emph{Say it:} \emph{vṛddhasya lalāṭe valiḥ asti.} — There is a wrinkle on the old man's forehead.
  \item \emph{Recall:} say the Sanskrit for length of life, then the Sanskrit for grey-haired, then say \emph{valiḥ} again
\end{itemize}

\begin{tcolorbox}[breakable,colback=teal!4,colframe=teal!35!black,title={Before you move on}]
What does \sk{वलिः} mean? (\textquotedblleft{}A wrinkle\textquotedblright{}.) Say it once more.
\end{tcolorbox}
\section[\texorpdfstring{\sk{रूपम्} (rūpam)}{रूपम् (rūpam)}]{\sk{रूपम्} (rūpam) — appearance; beauty}

\label{lesson:SA-C490-rupam}



\begin{quote}
Before the new one: say the Sanskrit for grey-haired, then the Sanskrit for a wrinkle.
\end{quote}

\subsection*{You'll want to know: \sk{रूपम्}}
\textbf{\sk{रूपम्}} — \emph{rūpam} — \textquotedblleft{}appearance; beauty\textquotedblright{}.

\begin{quote}
\sk{मीरायाः} \sk{रूपं} \sk{मातुः} \sk{इव} \sk{अस्ति।}
\end{quote}

\begin{quote}
\emph{mīrāyāḥ rūpaṃ mātuḥ iva asti.}
\end{quote}

\begin{quote}
Mira looks like her mother.
\end{quote}

A neuter noun.

\begin{cousinweb}
\textbf{\sk{रूपम्}} (\emph{rūpam}) MW suggests, with a \textquotedblleft{}perhaps\textquotedblright{}, a link with \emph{varpa}, \emph{varpas}, \textquotedblleft{}form\textquotedblright{}; it gives \textquotedblleft{}form, shape, figure\textquotedblright{} and \textquotedblleft{}loveliness, beauty\textquotedblright{}, and its origin is not settled (MW).
\end{cousinweb}

\begin{grammarlens}[title={What you've built}]
One more word for describing what you lived and felt.
\end{grammarlens}

\subsection*{Your turn}
\begin{itemize}
\raggedright
  \item \emph{Say it:} \emph{rūpam}
  \item \emph{Say it:} \emph{rūpam}, once more
  \item \emph{Say it:} \emph{mīrāyāḥ rūpaṃ mātuḥ iva asti.} — Mira looks like her mother.
  \item \emph{Recall:} say the Sanskrit for grey-haired, then the Sanskrit for a wrinkle, then say \emph{rūpam} again
\end{itemize}

\begin{tcolorbox}[breakable,colback=teal!4,colframe=teal!35!black,title={Before you move on}]
What does \sk{रूपम्} mean? (\textquotedblleft{}Appearance; beauty\textquotedblright{}.) Say it once more.
\end{tcolorbox}
